\section{A quantitative constraint from braid groups}\label{sec:braids}

Let $B_l$ denote the group on $s_1,\ldots,s_{l-1}$ with relations
\[
 s_i s_{i+1}s_i=s_{i+1}s_i s_{i+1},\qquad
 s_i s_j=s_j s_i\quad\text{if }|i-j|\ge2.
\]
Throughout this section the same notation may denote elements of another
group satisfying these relations.  No injectivity or cocompactness of
the resulting braid subgroup is needed for the length estimates.

\begin{proposition}\label{prop:braid-angle}
Suppose a group acting geometrically on a CAT(0) space contains elements
$s_1,\ldots,s_{39}$ satisfying the $B_{40}$ relations, with
$\ell(s_1^2)>0$.  Then
\begin{equation}\label{eq:braid-angle-bound}
 \left(\frac{\ell((s_1s_2)^3)}{3\ell(s_1^2)}\right)^2
 \ge\frac{143}{513}>\frac14,
\end{equation}
and in particular
\[
 \frac{\ell((s_1s_2)^3)}{3\ell(s_1^2)}
 =\frac{\ell((s_1s_2)^3)}{3\ell(s_2^2)}>\frac12.
\]
\end{proposition}

For $B_4$, Crisp and Paoluzzi compute translation components along an
axis of a central element under a minimal geometric action
\cite[Proposition~5]{CrispPaoluzzi2005}; minimality means that no proper
nonempty closed convex subset is invariant.

To prove Proposition~\ref{prop:braid-angle}, we compare central elements
of the initial braid subblocks using the centralizer characters of
Section~\ref{sec:lengths}.  For a general $B_l$ with $l\ge3$, define
\begin{equation}\label{eq:braid-centers-definition}
 L_1=1,\qquad L_j=s_{j-1}L_{j-1}s_{j-1},\qquad
 Z_r=L_1\cdots L_r,\qquad N_r=\binom r2.
\end{equation}
Here $2\le j\le l$ and $1\le r\le l$.

\begin{lemma}\label{lem:braid-centers}
The element $L_j$ commutes with $s_i$ for $i\le j-2$, and $Z_r$ commutes
with every $s_i$ for $i<r$.  In particular,
\[
 Z_2=s_1^2,\qquad Z_3=(s_1s_2)^3.
\]
The element $\delta=s_1\cdots s_{l-1}$ satisfies
\begin{equation}\label{eq:braid-shift}
 \delta s_i\delta^{-1}=s_{i+1}\qquad(1\le i\le l-2).
\end{equation}
In any geometric action of the ambient group,
\begin{equation}\label{eq:braid-central-pairing}
 \langle s_i^2,Z_r\rangle=\frac{\ell(Z_r)^2}{N_r}
 \qquad(2\le r\le l,\ 1\le i<r).
\end{equation}
\end{lemma}

\begin{proof}
Induct on $j$ for the assertion about $L_j$.  For $i\le j-3$, both
$s_{j-1}$ and $L_{j-1}$ commute with $s_i$.  For $i=j-2$, put
$a=s_{j-2}$, $b=s_{j-1}$, and $C=L_{j-2}$.  Then $bC=Cb$ and
$L_j=baCab$, so the braid relation gives
\[
 aL_j=abaCab=babCab=baCbab=baCaba=L_j a.
\]
This includes the case $j=3$, where $C=L_1=1$; the case $j=2$ is empty.
In particular, $L_j$ commutes with each earlier $L_i$, because the word
for $L_i$ uses only $s_1,\ldots,s_{i-1}$.

Next induct on $r$ for $Z_r$.  Both $Z_{r-1}$ and $L_r$ commute with
$s_i$ for $i\le r-2$.  The only remaining generator is $s=s_{r-1}$.
It commutes with $Z_{r-2}$.  With $L=L_{r-1}$ we have $L_r=sLs$,
and the already proved commutation of $L$ with $L_r$ says
$LsLs=sLsL$.  Therefore
\[
 s(LL_r)=sLsLs=LsLs^2=(LL_r)s,
\]
which proves the assertion.  The case $r=2$ is immediate.  Finally
\[
 Z_3=s_1^2s_2s_1^2s_2
     =s_1(s_1s_2s_1)s_1s_2
     =s_1(s_2s_1s_2)s_1s_2=(s_1s_2)^3.
\]

To verify \eqref{eq:braid-shift}, in $\delta s_i$ commute the last
$s_i$ left past $s_{i+2},\ldots,s_{l-1}$, replace
$s_i s_{i+1}s_i$ by $s_{i+1}s_i s_{i+1}$, and commute the new
$s_{i+1}$ left past $s_1,\ldots,s_{i-1}$.  The result is
$s_{i+1}\delta$.

For the length identity, apply the character
$h\mapsto\langle Z_r,h\rangle$ from
Lemma~\ref{lem:centralizer-character}.  All $s_i$, $i<r$, are in its
domain.  Applying it to each adjacent braid relation shows that their
values are equal; call the common value $\alpha$.  The positive word
for $L_j$ has $2(j-1)$ letters, so that for $Z_r$ it has $2N_r$ letters.
Thus
\[
 \ell(Z_r)^2=\langle Z_r,Z_r\rangle=2N_r\alpha,
 \qquad \langle s_i^2,Z_r\rangle=2\alpha,
\]
which proves \eqref{eq:braid-central-pairing}, including when
$\ell(Z_r)=0$.
\end{proof}

\begin{proof}[Proof of Proposition~\ref{prop:braid-angle}]
We retain $l=2k$ with $k\ge3$ until the final arithmetic, where $k=20$.
Put
$g=s_1^2$, $D=\ell(g)>0$, and
\[
 R_r=\frac{\ell(Z_r)^2}{N_r^2D^2}\quad(2\le r\le 2k),\qquad
 c_0=\frac{\langle s_1^2,s_3^2\rangle}{D^2}.
\]
We will express each normalized center length $R_r$ in terms of $R_3$
and $c_0$.  Positivity for the commuting odd-indexed squares and for two
disjoint block centers will then constrain these quantities and force
$R_3>1/4$.

All squared generators have length $D$, by
\eqref{eq:braid-shift}.  We claim that every pair with $j-i\ge2$ has
\begin{equation}\label{eq:nonadjacent-pairing}
 \langle s_i^2,s_j^2\rangle=c_0D^2.
\end{equation}
First fix $i=1$.  All $s_j$ with $j\ge3$ commute with $g$; their
adjacent braid relations, evaluated under the character of $g$, show
that $\langle g,s_j\rangle$ is independent of $j\ge3$.
This proves the claim for $i=1$.  For general $i,j$, conjugate the pair
$(s_1,s_{j-i+1})$ by $\delta^{i-1}$.  Every successive use of
\eqref{eq:braid-shift} has index at most $l-2$, so it sends the pair
to $(s_i,s_j)$.  Simultaneous conjugation invariance proves the claim.

By \eqref{eq:braid-central-pairing},
\begin{equation}\label{eq:normalized-center-pairing}
 \langle g,Z_r\rangle=N_rR_rD^2.
\end{equation}
We compute this character in a second way.  Since $Z_3=gL_3$,
\[
 \langle g,L_3\rangle=(3R_3-1)D^2.
\]
For $j>3$ the recursive definition gives the word
\[
 L_j=(s_{j-1}\cdots s_3)L_3(s_3\cdots s_{j-1}).
\]
Every displayed factor belongs to $C_G(g)$, and
$\langle g,s_i\rangle=c_0D^2/2$ for $i\ge3$.  Keeping $L_3$ as one
factor, the character therefore gives
\[
 \langle g,L_j\rangle
   =\bigl(3R_3-1+(j-3)c_0\bigr)D^2\qquad(j\ge3).
\]
In particular, this calculation does not expand the noncommuting
$s_2$-letters inside $L_3$.  Summing over the factors of $Z_r$, including
$L_2=g$, and using \eqref{eq:normalized-center-pairing}, yields
\begin{equation}\label{eq:braid-recursion}
 N_rR_r=1+(r-2)(3R_3-1)+\binom{r-2}{2}c_0
 \qquad(3\le r\le 2k).
\end{equation}

There are $k$ mutually commuting squares
$s_1^2,s_3^2,\ldots,s_{2k-1}^2$.  The squared length of their product,
computed by Lemma~\ref{lem:commuting-form}, is
\[
 kD^2+k(k-1)c_0D^2\ge0.
\]
Consequently
\begin{equation}\label{eq:lower-cross-pairing}
 c_0\ge-\frac1{k-1}.
\end{equation}

For the other constraint, compare $Z_k$ with
$Z'_k=\delta^k Z_k\delta^{-k}$.  Repeated application of
\eqref{eq:braid-shift} sends its generators $s_1,\ldots,s_{k-1}$ to
$s_{k+1},\ldots,s_{2k-1}$; the largest index to which the shift formula
is applied is $2k-2$.  Every generator in the first block commutes with
every generator in the second, so $Z_k$ and $Z'_k$ commute, and they have
equal length by conjugacy.

Write their positive words as $Z_k=x_1\cdots x_{2N_k}$ and
$Z'_k=y_1\cdots y_{2N_k}$.  Use first the character of $Z'_k$, since
each $x_p$ commutes with $Z'_k$, and then the character of each $x_p$,
since each $y_q$ commutes with $x_p$.  These are two separate valid
character expansions, giving
\begin{align*}
 \langle Z_k,Z'_k\rangle
 &=\sum_{p=1}^{2N_k}\langle x_p,Z'_k\rangle
   =\sum_{p,q=1}^{2N_k}\langle x_p,y_q\rangle\\
 &= (2N_k)^2\frac{c_0D^2}{4}=N_k^2c_0D^2.
\end{align*}
Here the factor $1/4$ follows by bilinearity from
\eqref{eq:nonadjacent-pairing} applied to $x_p^2,y_q^2$.
The nonnegative squared length of $Z_k(Z'_k)^{-1}$ is thus
$2N_k^2D^2(R_k-c_0)$, proving
\begin{equation}\label{eq:disjoint-centers}
 R_k\ge c_0.
\end{equation}

Put $r=k$ in \eqref{eq:braid-recursion} and use
\eqref{eq:disjoint-centers}.  Since
$N_k-\binom{k-2}{2}=2k-3$, we obtain
\begin{align*}
 3R_3
 &\ge 1-\frac1{k-2}+\frac{2k-3}{k-2}c_0\\
 &\ge 1-\frac{3k-4}{(k-1)(k-2)},
\end{align*}
where the second inequality is \eqref{eq:lower-cross-pairing}.
At $k=20$, this says
\[
 3R_3\ge 1-\frac{56}{342}=\frac{143}{171},\qquad
 R_3\ge\frac{143}{513}>\frac14.
\]
Finally $Z_3=(s_1s_2)^3$, $N_3=3$, and
$\ell(s_2^2)=\ell(s_1^2)=D$.  These identities give the claimed ratios.
\end{proof}
